% !PRINTED 476
\weberpage{476}
\textbf{§\,129. \gesperrt{Application of the Transformation of the Second Order to the
Calculation of Class Invariants.}}

The transformation of the second order can be applied, repeating it
several times if necessary, to all those discriminants
$-4m$ for which $y^2 + m x^2$ is a power of 2, thus
for example to $m = 7$, $15$, $23$, $31$.

These calculations are, however, laborious, and we shall find
simpler ways of reaching the goal in these cases.
The transformation of the second order renders better service in
finding, from a known class invariant, a new one belonging to
a discriminant which is four times the former.

The following formulae lead to this:
According to § 34:
\[
f_1(\omega)^8 + f_2(\omega)^8 = f(\omega)^8,
\]
\[
f_1(\omega)^8 f_2(\omega)^8 = \frac{16}{f(\omega)^8},
\]
from which
\[
f_1(\omega)^8 - f_2(\omega)^8 = \frac{\sqrt{f(\omega)^{24} - 64}}{f(\omega)^4}.
\]
The sign of the root changes only for $f(\omega)^{24} = 64$,
that is, for $\omega = i$, and so long as $-i\omega$ is real and greater than 1
it is to be taken positive, since the left-hand side becomes
positively infinite for a vanishing $q$. From this it follows:
\[
2 f_2(\omega)^8 = \frac{f(\omega)^{12} - \sqrt{f(\omega)^{24} - 64}}{f(\omega)^4},
\]
or
\[
f_2(\omega)^8 f(\omega)^4 \left[f(\omega)^{12} + \sqrt{f(\omega)^{24} - 64}\right] = 32,
\]
and in the same way:
\[
f_2(\omega)^8 f_1(\omega)^4 \left[f_1(\omega)^{12} + \sqrt{f_1(\omega)^{24} + 64}\right] = 32.
\]
Further, according to § 34, (16):
\[
f_1(2\omega)\,f_2(\omega) = \sqrt{2},
\]
from which
\weq{1}{\begin{aligned}
2 f_1(2\omega)^8 &= f(\omega)^4\left[f(\omega)^{12} + \sqrt{f(\omega)^{24} - 64}\right],\\
&= f_1(\omega)^4\left[f_1(\omega)^{12} + \sqrt{f_1(\omega)^{24} + 64}\right].
\end{aligned}}
These formulae can serve, when $f(\omega)$ or $f_1(\omega)$ is
known, to calculate $f_1(2\omega)$, and so $f_1(\sqrt{-4m})$ from $f(\sqrt{-m})$
%CONT>>
